\section*{September 4, 2026: Amenity controls leave the price estimates largely unchanged}
\textit{Agent-drafted research entry; not yet reviewed by Jacob.}

Allowing prices to vary with proximity to the lake, parks, and rail stations
in each half-year leaves the dollar estimate near \$58,000 and modestly
reduces the log estimate. The comparison uses sales within a quarter mile of
the selected 30 closed and 49 still-open schools. Both models account for
housing characteristics and persistent price differences across school areas;
only the amenity controls differ.

Each transaction receives distances to the Lake Michigan shoreline, the
nearest park in the August 2012 inventory, and the nearest station in an
archived 2012 CTA network. Each distance enters linearly in miles with a
separate coefficient in every calendar half-year. These relationships are
common across treated and control homes. The event study therefore estimates
66 amenity slopes; the pooled DiD, which excludes both halves of 2013,
estimates 60. This lets, for example, proximity to the lake predict a different rate of
price growth over time.

\begin{center}
\small
\input{../output/amenity_table.tex}
\end{center}
{\footnotesize Both specifications include housing characteristics and
school-site and calendar-half-year fixed effects. The pooled DiD compares
2008--2012 with 2014--2018. Log estimates use
$100[\exp(\hat\beta)-1]$; dollar estimates use 2022 purchasing power.
Pre-period tests come from semiannual event studies normalized to 2012 H2.}

The pooled comparison uses 9,912 sales, and the event studies use 10,921,
across 29 treated and 48 control sites. Adding amenities removes no
transactions. The dollar interval is narrower and now excludes zero at the
5\% level, but the estimate itself barely changes. The log estimate remains
imprecise. This particular adjustment does not explain the difference
between dollar and log estimates or establish parallel trends.

The park boundaries describe August 2012. The CTA layer is
an Esri-hosted preservation of City of Chicago data whose metadata identifies
a March 2013 download of a 2012 network; it contains 145 stations. This
vintage relies on the preservation metadata, not an independently recovered
historical transit feed. The lake feature carries a 2005 edit date. Baseline
amenity locations stay fixed across sale years, including 2008--2011; the
measures are not annual inventories or measures of service quality.

The distance task preserves all 167,468 clean transactions. One lacks a
centroid and retains missing distances; it is outside the preferred
estimation sample. The distance-data report records missing values and distance
distributions. One self-intersecting park
polygon, Douglas (Stephen), is repaired in memory without dropping it.
The 2013 Red Line South reconstruction is not separately modeled, and the
specification includes neither bus nor Metra access.

\clearpage
\begin{center}
\includegraphics[width=\textwidth]{../input/event_studies_amenities.png}
\end{center}
{\footnotesize Top panels include housing characteristics; bottom panels add
amenity distances interacted with half-year. Vertical axes match within each
outcome. The hollow point is the normalized 2012 H2 reference, shaded 2013
is the transition year, and intervals are pointwise 95\% intervals clustered
by school site. A post-period log coefficient near zero can still represent
an increase from the earlier negative pre-period coefficients.}

\textit{Reproduction.} Run \texttt{make} in the
\texttt{tasks/audits/home\_prices\_twfe/code} directory; upstream Make rules
build the new amenity inputs. Build this entry with \texttt{make} in
\texttt{logbook/code}. Source URLs and geographic-vintage qualifications are
in the \texttt{download\_baseline\_amenities} audit README. The
\texttt{home\_amenity\_distances} audit owns the transaction-distance CSV
and its data report. Code revision: \texttt{c200c31} plus current uncommitted
changes; geography downloaded September 4, 2026. Housing inputs are the
previously inspected September 4 sample. Independent dummy-variable OLS
checks and direct distance comparisons accompany verification of this run.
